

%%%%%%%%%%%%%%%%%%%%%%%%

%

% $Autor: Manoj Kumar Prabhakaran $

% $Datum: 2025-03-22 $

% $Pfad: GitHub/ML24-EX03-Pico4MLPro/Pico4MLPro/Contents/General/ILI9341Display.tex $

% $Version: 1 $

%

%%%%%%%%%%%%%%%%%%%%%%%%












\chapter{ILI9341 Display}





\section{Description}


The ILI9341 TFT display is a high-resolution color display module commonly integrated into embedded systems, particularly with microcontrollers like the Raspberry Pi Pico4MLPro. This display provides vibrant color output with touchscreen capabilities, making it suitable for interactive embedded applications requiring graphical user interfaces \cite{Adafruit:2025}.





\section{Hardware Interfaces}





\subsection{Pin Configuration}


The display interfaces with the Raspberry Pi Pico4MLPro through the following connections \cite{MichielvanderWulp:2024}:





\begin{table}[h]
	
	
	\centering
	
	
	\begin{tabular}{|l|l|}
		
		
		\hline
		
		
		\textbf{Display Pin} & \textbf{Pico4MLPro GPIO} \\
		
		
		\hline
		
		
		VCC & 3.3V/5V \\
		
		
		GND & GND \\
		
		
		CS & GP5 \\
		
		
		DC & GP6 \\
		
		
		RESET & GP7 \\
		
		
		MOSI & GP3 \\
		
		
		SCK & GP2 \\
		
		
		\hline
		
		
	\end{tabular}
	
	
	\caption{ILI9341 to Pico4MLPro Connection}
	
	
\end{table}





\subsection{Interface Protocols}


The display supports multiple interface protocols \cite{MicroRobotics:2024}:


\begin{itemize}
	
	
	\item 4-wire SPI (Serial Peripheral Interface)
	
	
	\item 8-bit parallel
	
	
	\item 16-bit parallel
	
	
\end{itemize}





For the Pico4MLPro implementation, the 4-wire SPI interface is utilized due to its minimal pin requirements and adequate performance characteristics.





\section{Specifications}





\subsection{Display Characteristics}


\begin{itemize}
	
	
	\item \textbf{Resolution}: 320 × 240 pixels \cite{PCBOnline:2024}
	
	
	\item \textbf{Color Depth}: 16-bit (65,536 colors) or 18-bit (262,144 colors) \cite{Adafruit:2025}
	
	
	\item \textbf{Controller}: ILI9341 display driver
	
	
	\item \textbf{Touch Controller}: XPT2046 (when touchscreen functionality is included) \cite{MichielvanderWulp:2024}
	
	
\end{itemize}





\subsection{Electrical Specifications}


\begin{itemize}
	
	
	\item \textbf{Operating Voltage}: 3.3V logic level required regardless of power supply voltage \cite{MichielvanderWulp:2024}
	
	
	\item \textbf{Power Consumption}: 80mA backlight current, requiring separate power regulation for battery-operated implementations \cite{MichielvanderWulp:2024}
	
	
\end{itemize}





\section{Constraints}





\subsection{Performance Limitations}


\begin{itemize}
	
	
	\item \textbf{SPI Speed Tradeoffs}: Maximum 62.5 MHz SPI clock induces approximately 5\% CPU load during DMA transfers \cite{ArduinoForum:2025}
	
	
	\item \textbf{Memory Requirements}: Full-screen 16-bit color buffer requires 153.6 KB of RAM, exceeding some microcontrollers' capabilities \cite{HJWWalters:2025}
	
	
\end{itemize}





\subsection{Implementation Considerations}


\begin{itemize}
	
	
	\item Logic levels must remain at 3.3V regardless of the power supply voltage used for the display \cite{MichielvanderWulp:2024}
	
	
	\item Separate power regulation is required for battery-operated projects due to the backlight current draw \cite{MichielvanderWulp:2024}
	
	
\end{itemize}





\section{Dimensions}


The ILI9341 TFT display module standard dimensions are compatible with the Pico4MLPro form factor, though specific dimensions vary by manufacturer implementation.





\section{OS/Software Interfaces/Protocols}





\subsection{Library Support}


Two primary software libraries are recommended for interfacing with the ILI9341 display:





\begin{enumerate}
	
	
	\item \textbf{TFT\_eSPI Library} - Optimized for Raspberry Pi Pico with DMA support \cite{ArduinoForum:2025}
	
	
	\begin{itemize}
		
		
		\item Achieves up to 43 FPS at 62.5 MHz SPI clock
		
		
		\item Compatible with Arduino IDE and Pico SDK
		
		
	\end{itemize}
	
	
	
	
	
	\item \textbf{CircuitPython Library} - Adafruit's hardware-accelerated rendering solution \cite{PCBWay:2025}
	
	
	\begin{itemize}
		
		
		\item Simpler implementation for Python-based projects
		
		
		\item Built-in graphics primitives and text rendering
		
		
	\end{itemize}
	
	
\end{enumerate}





\subsection{SPI Protocol Configuration}


\begin{itemize}
	
	
	\item \textbf{Clock Speed}: Up to 62.5 MHz
	
	
	\item \textbf{Mode}: SPI Mode 0 (CPOL=0, CPHA=0)
	
	
	\item \textbf{Bit Order}: MSB First
	
	
\end{itemize}





\subsection{Touch Interface Protocol}


The resistive touch layer uses the XPT2046 controller with:


\begin{itemize}
	
	
	\item \textbf{Sampling Rate}: 125 Hz
	
	
	\item \textbf{Interface}: Separate SPI channel or shared with display
	
	
	\item \textbf{Calibration}: Requires 4-point adjustment for accurate touch mapping
	
	
\end{itemize}





\section{Data}





\subsection{Quality}


The ILI9341 TFT display provides professional-grade display quality suitable for embedded systems applications \cite{PCBOnline:2024}:


\begin{itemize}
	
	
	\item High contrast ratio
	
	
	\item Wide viewing angle
	
	
	\item Fast response time for smooth animations and UI interactions
	
	
\end{itemize}





\subsection{Quantity}


Memory utilization metrics for display operations \cite{HJWWalters:2025}:


\begin{itemize}
	
	
	\item Full-screen 16-bit color buffer: 153.6 KB
	
	
	\item Pico4MLPro's 264 KB SRAM enables complex GUI implementations without external memory \cite{PCBWay:2025}
	
	
\end{itemize}





\subsection{Types}


The display supports various data formats:


\begin{itemize}
	
	
	\item RGB565 (16-bit color)
	
	
	\item RGB666 (18-bit color)
	
	
	\item Various text and font rendering options through libraries
	
	
\end{itemize}





\subsection{Structure}


Data transmission follows standard SPI protocol structure with:


\begin{itemize}
	
	
	\item Command/Data selection via DC pin
	
	
	\item 8-bit command bytes
	
	
	\item Variable-length data parameters
	
	
	\item Hardware-based chip selection via CS pin
	
	
\end{itemize}





\section{Code with Simple Example}

\subsection{TFT\_eSPI Library Example (C++)}

{
	\captionof{code}{Example ILI9341 Display code implementation}\label{ILI9341Display:exampleard}
	\ArduinoExternal{}{../../Code/Pico4MLPro/ILI9341Display/ILI9341Display.ino}
}

\subsection{CircuitPython Example}

{
	\begin{lstlisting}[language=MyPython]
import board
import busio
import digitalio
import adafruit_ili9341
		
# Configure SPI
spi = busio.SPI(clock=board.GP2, MOSI=board.GP3)
		
# Configure chip select and DC pins
cs_pin = digitalio.DigitalInOut(board.GP5)
dc_pin = digitalio.DigitalInOut(board.GP6)
		
# Initialize display
display = adafruit_ili9341.ILI9341(
               spi,
               cs=cs_pin,
               dc=dc_pin,
               width=320,
               height=240
           )
		
# Fill screen with color
display.fill(0x0000FF)  # Blue
		
# Draw text
display.text("Hello Pico4MLPro", 10, 10, 0xFFFFFF)
  \end{lstlisting}
	
	
	
	
	
	
	
	
	
	
	
	
	
	
	
	
	\section{Performance Metrics}
	
	
	\begin{itemize}
		
		
		\item \textbf{Frame Rate}: 43 FPS achievable with DMA \cite{ArduinoForum:2025}
		
		
		\item \textbf{Memory Requirements}:
		
		
		\begin{itemize}
			
			
			\item Full-screen 16-bit color: 153.6 KB (exceeds Arduino Uno's 32 KB RAM) \cite{HJWWalters:2025}
			
			
			\item Pico4MLPro's 264 KB SRAM enables complex GUI implementations \cite{PCBWay:2025}
			
			
		\end{itemize}
		
		
	\end{itemize}